% Algebra 2 Lesson 3-7 fast-track packet: Transformations of Polynomial Functions
% (enVision A2 2024 TE, Example 3 only: brief lesson; Examples 1, 2, 4 and 5 cut, fast-lesson spec R1-R3).
% Quiz items fed: 2 (parent of -x^4 - 1 and the order of its transformations) and 3 (how 3x^3 + 6 differs
% from x^3) by Example 3; 4 (equation of a quartic from its vertex, bonus) by Example 3 read backwards.
% DOK 3: Practice 12 (TE files it under Example 2), as the Extension.
% Shared preamble for the "fast" (hacked-down) Algebra 2 lesson packets, Oct 2026.
% Each packet: % Shared preamble for the "fast" (hacked-down) Algebra 2 lesson packets, Oct 2026.
% Each packet: % Shared preamble for the "fast" (hacked-down) Algebra 2 lesson packets, Oct 2026.
% Each packet: \input{fast-packet-template} then \begin{document} ... \end{document}.
\documentclass[12pt]{article}
\usepackage[margin=0.7in]{geometry}
\usepackage{amsmath,amssymb}
\usepackage{booktabs,tabularx,array}
\usepackage{enumitem}
\usepackage[T1]{fontenc}
\usepackage[utf8]{inputenc}
\usepackage{lmodern}
\usepackage{tikz}
\usepackage{pgfplots}
\pgfplotsset{compat=1.18}
\usepackage{xcolor}
\setlist{nosep,leftmargin=1.4em}
\setlength{\parindent}{0pt}
\setlength{\parskip}{6pt}
\renewcommand{\arraystretch}{1.35}
\newcolumntype{Y}{>{\raggedright\arraybackslash}X}
\pagestyle{empty}

% blanks and workspace
\newcommand{\blank}[1][2.6cm]{\rule[-2pt]{#1}{0.5pt}}
\newcommand{\lines}[1][1]{\par\vspace{2pt}\foreach \i in {1,...,#1}{\rule{0pt}{22pt}\rule{\linewidth}{0.4pt}\par}}
\newcommand{\wspart}[2]{\vspace{4pt}{\large\bfseries #1: #2}\par\vspace{-6pt}\rule{\textwidth}{0.4pt}\par}

% header: \packethead{Lesson 1-3}{Piecewise-Defined Functions}
\newcommand{\packethead}[2]{%
{\LARGE\bfseries #2}\hfill Name \blank[5cm]\\[2pt]
Algebra 2 \quad #1 (fast track) \hfill Date \blank[2.5cm] \quad Period \blank[1.2cm]\par\vspace{2pt}}

% boxed goals: \goals{learning goal}{language goal}
\newcommand{\goals}[2]{%
\fbox{\parbox{\dimexpr\textwidth-2\fboxsep-2\fboxrule}{%
\textbf{Learning goal:} #1\\[2pt]
\textbf{Language goal:} #2}}\par}

% quiz target box: \quiztarget{text}
\newcommand{\quiztarget}[1]{%
\vspace{4pt}\colorbox{gray!12}{\parbox{\dimexpr\textwidth-2\fboxsep}{\textbf{On the quiz:} #1}}\par}

% IXL table row helper: code & skill & which quiz question it feeds
\newcommand{\ixlrow}[3]{\texttt{#1} & #2 & #3 \\}

% one grid style for every graph; override xmin/xmax/... per plot
\pgfplotsset{sheetgrid/.style={
  axis lines=middle, grid=both, grid style={gray!45},
  xmin=-10.5, xmax=10.5, ymin=-10.5, ymax=10.5,
  xtick={-10,-8,...,10}, ytick={-10,-8,...,10},
  x=0.28cm, y=0.28cm,
  ticklabel style={font=\scriptsize, fill=white, inner sep=1pt},
  xlabel={$x$}, ylabel={$y$},
  axis line style={thick}, clip=true}}
% unit grid for small windows
\pgfplotsset{unitgrid/.style={sheetgrid,
  xmin=-5.5, xmax=5.5, ymin=-5.5, ymax=5.5,
  xtick={-5,...,5}, ytick={-5,...,5}, x=0.5cm, y=0.5cm}}
% open and closed endpoints
\tikzset{closed/.style={circle, fill, inner sep=1.6pt},
         open/.style={circle, draw, fill=white, inner sep=1.4pt}}
 then \begin{document} ... \end{document}.
\documentclass[12pt]{article}
\usepackage[margin=0.7in]{geometry}
\usepackage{amsmath,amssymb}
\usepackage{booktabs,tabularx,array}
\usepackage{enumitem}
\usepackage[T1]{fontenc}
\usepackage[utf8]{inputenc}
\usepackage{lmodern}
\usepackage{tikz}
\usepackage{pgfplots}
\pgfplotsset{compat=1.18}
\usepackage{xcolor}
\setlist{nosep,leftmargin=1.4em}
\setlength{\parindent}{0pt}
\setlength{\parskip}{6pt}
\renewcommand{\arraystretch}{1.35}
\newcolumntype{Y}{>{\raggedright\arraybackslash}X}
\pagestyle{empty}

% blanks and workspace
\newcommand{\blank}[1][2.6cm]{\rule[-2pt]{#1}{0.5pt}}
\newcommand{\lines}[1][1]{\par\vspace{2pt}\foreach \i in {1,...,#1}{\rule{0pt}{22pt}\rule{\linewidth}{0.4pt}\par}}
\newcommand{\wspart}[2]{\vspace{4pt}{\large\bfseries #1: #2}\par\vspace{-6pt}\rule{\textwidth}{0.4pt}\par}

% header: \packethead{Lesson 1-3}{Piecewise-Defined Functions}
\newcommand{\packethead}[2]{%
{\LARGE\bfseries #2}\hfill Name \blank[5cm]\\[2pt]
Algebra 2 \quad #1 (fast track) \hfill Date \blank[2.5cm] \quad Period \blank[1.2cm]\par\vspace{2pt}}

% boxed goals: \goals{learning goal}{language goal}
\newcommand{\goals}[2]{%
\fbox{\parbox{\dimexpr\textwidth-2\fboxsep-2\fboxrule}{%
\textbf{Learning goal:} #1\\[2pt]
\textbf{Language goal:} #2}}\par}

% quiz target box: \quiztarget{text}
\newcommand{\quiztarget}[1]{%
\vspace{4pt}\colorbox{gray!12}{\parbox{\dimexpr\textwidth-2\fboxsep}{\textbf{On the quiz:} #1}}\par}

% IXL table row helper: code & skill & which quiz question it feeds
\newcommand{\ixlrow}[3]{\texttt{#1} & #2 & #3 \\}

% one grid style for every graph; override xmin/xmax/... per plot
\pgfplotsset{sheetgrid/.style={
  axis lines=middle, grid=both, grid style={gray!45},
  xmin=-10.5, xmax=10.5, ymin=-10.5, ymax=10.5,
  xtick={-10,-8,...,10}, ytick={-10,-8,...,10},
  x=0.28cm, y=0.28cm,
  ticklabel style={font=\scriptsize, fill=white, inner sep=1pt},
  xlabel={$x$}, ylabel={$y$},
  axis line style={thick}, clip=true}}
% unit grid for small windows
\pgfplotsset{unitgrid/.style={sheetgrid,
  xmin=-5.5, xmax=5.5, ymin=-5.5, ymax=5.5,
  xtick={-5,...,5}, ytick={-5,...,5}, x=0.5cm, y=0.5cm}}
% open and closed endpoints
\tikzset{closed/.style={circle, fill, inner sep=1.6pt},
         open/.style={circle, draw, fill=white, inner sep=1.4pt}}
 then \begin{document} ... \end{document}.
\documentclass[12pt]{article}
\usepackage[margin=0.7in]{geometry}
\usepackage{amsmath,amssymb}
\usepackage{booktabs,tabularx,array}
\usepackage{enumitem}
\usepackage[T1]{fontenc}
\usepackage[utf8]{inputenc}
\usepackage{lmodern}
\usepackage{tikz}
\usepackage{pgfplots}
\pgfplotsset{compat=1.18}
\usepackage{xcolor}
\setlist{nosep,leftmargin=1.4em}
\setlength{\parindent}{0pt}
\setlength{\parskip}{6pt}
\renewcommand{\arraystretch}{1.35}
\newcolumntype{Y}{>{\raggedright\arraybackslash}X}
\pagestyle{empty}

% blanks and workspace
\newcommand{\blank}[1][2.6cm]{\rule[-2pt]{#1}{0.5pt}}
\newcommand{\lines}[1][1]{\par\vspace{2pt}\foreach \i in {1,...,#1}{\rule{0pt}{22pt}\rule{\linewidth}{0.4pt}\par}}
\newcommand{\wspart}[2]{\vspace{4pt}{\large\bfseries #1: #2}\par\vspace{-6pt}\rule{\textwidth}{0.4pt}\par}

% header: \packethead{Lesson 1-3}{Piecewise-Defined Functions}
\newcommand{\packethead}[2]{%
{\LARGE\bfseries #2}\hfill Name \blank[5cm]\\[2pt]
Algebra 2 \quad #1 (fast track) \hfill Date \blank[2.5cm] \quad Period \blank[1.2cm]\par\vspace{2pt}}

% boxed goals: \goals{learning goal}{language goal}
\newcommand{\goals}[2]{%
\fbox{\parbox{\dimexpr\textwidth-2\fboxsep-2\fboxrule}{%
\textbf{Learning goal:} #1\\[2pt]
\textbf{Language goal:} #2}}\par}

% quiz target box: \quiztarget{text}
\newcommand{\quiztarget}[1]{%
\vspace{4pt}\colorbox{gray!12}{\parbox{\dimexpr\textwidth-2\fboxsep}{\textbf{On the quiz:} #1}}\par}

% IXL table row helper: code & skill & which quiz question it feeds
\newcommand{\ixlrow}[3]{\texttt{#1} & #2 & #3 \\}

% one grid style for every graph; override xmin/xmax/... per plot
\pgfplotsset{sheetgrid/.style={
  axis lines=middle, grid=both, grid style={gray!45},
  xmin=-10.5, xmax=10.5, ymin=-10.5, ymax=10.5,
  xtick={-10,-8,...,10}, ytick={-10,-8,...,10},
  x=0.28cm, y=0.28cm,
  ticklabel style={font=\scriptsize, fill=white, inner sep=1pt},
  xlabel={$x$}, ylabel={$y$},
  axis line style={thick}, clip=true}}
% unit grid for small windows
\pgfplotsset{unitgrid/.style={sheetgrid,
  xmin=-5.5, xmax=5.5, ymin=-5.5, ymax=5.5,
  xtick={-5,...,5}, ytick={-5,...,5}, x=0.5cm, y=0.5cm}}
% open and closed endpoints
\tikzset{closed/.style={circle, fill, inner sep=1.6pt},
         open/.style={circle, draw, fill=white, inner sep=1.4pt}}

\begin{document}

\packethead{Lesson 3-7}{Transformations of Polynomial Functions}
\goals{I can describe how the graph of a cubic or quartic function is stretched, reflected, and translated from its parent function.}{I can explain each change to a graph using the words parent function, leading coefficient, stretch, reflection, and translation.}

\quiztarget{(1) Even, odd or neither from a graph (full Savvas quiz only). (2) The parent function and the ordered transformations of a quartic with a negative leading coefficient (Example 3, Practice). (3) Every true statement about a stretched and translated cubic (Example 3, Try It 3). (4, bonus) A quartic's equation from its turning point (Example 3 backwards). (5) A cube's volume in new units (full Savvas quiz only).}
\topictarget{3}{Question 14: how $f(x)=2x^4-3$ differs from its parent: select all (Ex.\ 3).}

{\raggedright\wspart{Example 3}{Graph Transformations of Cubic and Quartic Parent Functions}}
How do transformed graphs compare to the graph of the parent function? The \textbf{parent functions} are $y=x^3$ (cubic) and $y=x^4$ (quartic). In $g(x)=a(x-h)^n+k$, the \textbf{leading coefficient} $a$ stretches the graph vertically when $|a|>1$, compresses it when $0<|a|<1$, and also reflects it across the $x$-axis when $a<0$; $h$ translates it horizontally (right when $h>0$) and $k$ vertically (up when $k>0$).

\begin{minipage}[t]{0.48\textwidth}
\vspace{0pt}
\textbf{A.} $g(x)=3x^4-18$
\begin{enumerate}[label=\textbf{\arabic*.}, itemsep=0pt, topsep=0pt]
\item The parent function is $f(x)=x^4$.
\item The leading coefficient, 3, \textbf{stretches} the graph vertically, making it narrower than the graph of the parent function.
\item Subtracting 18 \textbf{translates} the graph down 18 units.
\end{enumerate}
\vspace{3pt}\small\renewcommand{\arraystretch}{1.1}
\begin{tabular}{|c|c|c|}
\hline
on $f=x^4$ & $\times 3$ & $-18$: on $g$ \\ \hline
$(0,0)$ & $(0,0)$ & $(0,-18)$ \\
$(\pm1,1)$ & $(\pm1,3)$ & $(\pm1,-15)$ \\
$(\pm2,16)$ & $(\pm2,48)$ & $(\pm2,30)$ \\ \hline
\end{tabular}
\normalsize\renewcommand{\arraystretch}{1.35}

\vspace{2pt}\centering
\begin{tikzpicture}
\begin{axis}[sheetgrid, xmin=-2.6, xmax=2.6, ymin=-40, ymax=60,
  xtick={-2.5,-2,...,2.5}, xticklabels={,$-2$,,$-1$,,,,$1$,,$2$,},
  ytick={-40,-30,...,60}, yticklabels={,,,,,,$20$,,$40$,,},
  x=0.9cm, y=0.03cm]
\addplot[very thick, blue, domain=-2.47:2.47, samples=80, <->] {x^4};
\addplot[very thick, red, domain=-2.236:2.236, samples=80, <->] {3*x^4-18};
\node[closed, red] at (axis cs:0,-18) {};
\node[closed, red] at (axis cs:-2,30) {};
\node[closed, red] at (axis cs:2,30) {};
\node[font=\scriptsize, fill=white, inner sep=1pt, anchor=north east] at (axis cs:0,0) {$O$};
\node[font=\scriptsize, fill=white, inner sep=1pt, anchor=north west] at (axis cs:0.1,-21) {$(0,-18)$};
\node[font=\scriptsize, fill=white, inner sep=1pt, anchor=north west] at (axis cs:-1.82,28) {$(-2,30)$};
\node[font=\scriptsize, fill=white, inner sep=1pt, anchor=north east] at (axis cs:1.82,28) {$(2,30)$};
\node[font=\scriptsize, blue, fill=white, inner sep=1pt] at (axis cs:-1.4,12) {$f$};
\node[font=\scriptsize, red, fill=white, inner sep=1pt] at (axis cs:2.35,50) {$g$};
\end{axis}
\end{tikzpicture}
\end{minipage}\hfill
\begin{minipage}[t]{0.49\textwidth}
\vspace{0pt}
\textbf{B.} $h(x)=5(x-2)^3+4$
\begin{enumerate}[label=\textbf{\arabic*.}, itemsep=0pt, topsep=0pt]
\item The parent function is $f(x)=x^3$.
\item Subtracting 2 (before calculating the cube) \textbf{translates} the graph right 2 units.
\item Multiplying by 5 \textbf{stretches} the translated graph vertically, making it narrower.
\item Adding 4 translates the stretched graph up 4 units.
\end{enumerate}
\vspace{3pt}\small\renewcommand{\arraystretch}{1.1}
\begin{tabular}{|c|c|c|c|}
\hline
on $f=x^3$ & right 2 & $\times5$ & $+4$: on $h$ \\ \hline
$(-1,-1)$ & $(1,-1)$ & $(1,-5)$ & $(1,-1)$ \\
$(0,0)$ & $(2,0)$ & $(2,0)$ & $(2,4)$ \\
$(1,1)$ & $(3,1)$ & $(3,5)$ & $(3,9)$ \\ \hline
\end{tabular}
\normalsize\renewcommand{\arraystretch}{1.35}

\vspace{2pt}\centering
\begin{tikzpicture}
\begin{axis}[sheetgrid, xmin=-4, xmax=6, ymin=-20, ymax=30,
  xtick={-4,...,6}, xticklabels={,,$-2$,,,,$2$,,$4$,,},
  ytick={-20,-15,...,30}, yticklabels={,,$-10$,,,,$10$,,$20$,,},
  x=0.45cm, y=0.054cm]
\addplot[very thick, blue, domain=-2.65:3.05, samples=80, <->] {x^3};
\addplot[very thick, red, domain=0.35:3.70, samples=80, <->] {5*(x-2)^3+4};
\node[closed, red] at (axis cs:2,4) {};
\node[closed, red] at (axis cs:1,-1) {};
\node[closed, red] at (axis cs:3,9) {};
\node[font=\scriptsize, fill=white, inner sep=1pt, anchor=south east] at (axis cs:0,0) {$O$};
\node[font=\scriptsize, fill=white, inner sep=1pt, anchor=north west] at (axis cs:2.25,-4.2) {$(2,4)$};
\node[font=\scriptsize, blue, fill=white, inner sep=1pt] at (axis cs:-2.9,-14) {$f$};
\node[font=\scriptsize, red, fill=white, inner sep=1pt] at (axis cs:4.5,25) {$h$};
\end{axis}
\end{tikzpicture}
\end{minipage}

\vspace{2pt}
\small\textbf{Check:} $g(2)=3(16)-18=30$, $h(3)=5+4=9$ \checkmark\quad \textbf{Order:} stretch or reflect \emph{first}, then translate vertically.\normalsize

\textbf{Say it:} The \textbf{parent function} of $h$ is \blank[1.6cm]. The \textbf{leading coefficient} \blank[1cm] \blank[2.6cm] the graph vertically; the \textbf{translation} is \blank[1.5cm] right and \blank[1.5cm] up. No \textbf{reflection}: the leading coefficient is \blank[1.8cm].

\newpage
\wspart{Try It 3}{Compare to the parent function}
How does the graph of the function $g(x)=2x^3-5$ differ from the graph of its parent function?

\begin{minipage}[t]{0.50\textwidth}
\vspace{0pt}\raggedright
Parent function: $y=\blank[2cm]$

\vspace{4pt}
The leading coefficient, 2, \blank[4.2cm]

\vspace{4pt}
\blank[\linewidth]

\vspace{4pt}
Subtracting 5 \blank[5cm]

\vspace{4pt}
\blank[\linewidth]

\vspace{4pt}
Check one point: $(1,1)$ on $y=x^3$ goes to \blank[2.4cm] on $g$.

\vspace{4pt}
\small (Optional) Sketch the parent function and $g$ on the grid.\normalsize
\end{minipage}\hfill
\begin{minipage}[t]{0.46\textwidth}
\vspace{0pt}
\begin{center}
\begin{tikzpicture}
\begin{axis}[sheetgrid, xmin=-3.5, xmax=3.5, ymin=-25, ymax=25,
  xtick={-3,...,3}, ytick={-25,-20,...,25}, yticklabels={,$-20$,,$-10$,,,,$10$,,$20$,},
  x=0.85cm, y=0.12cm]
\end{axis}
\end{tikzpicture}
\end{center}
\end{minipage}

% Practice under R5: quiz item 2 (f(x) = -x^4 - 1: parent x^4, reflect over the x-axis, then translate down 1)
% has a reflection, and neither Example 3 nor Try It 3 does. TE Practice 18 (g(x) = -x^4 - 8, SEE EXAMPLE 3) is its
% structural parallel with different numbers; TE Practice 8 (Terrence, DOK 2, Example 3) is the reflection-order error.
\wspart{Practice}{When the leading coefficient is negative}
A \textbf{negative} leading coefficient adds a \textbf{reflection} across the $x$-axis: the ends of the graph point the opposite way from the parent function's ends.

\begin{minipage}[t]{0.55\textwidth}
\vspace{0pt}
\textbf{18.} How do the graphs of transformed functions compare to the graph of the parent function?
\[g(x)=-x^4-8\]
Parent function: $y=\blank[2cm]$

The leading coefficient of $-1$ \blank[4cm]

\blank[\linewidth]

Subtracting 8 \blank[5.4cm]

\vspace{2pt}
Which comes first, the reflection or the translation? Why?
\lines[2]
\end{minipage}\hfill
\begin{minipage}[t]{0.42\textwidth}
\vspace{0pt}
\begin{center}
\begin{tikzpicture}
\begin{axis}[sheetgrid, xmin=-5, xmax=3, ymin=-10.4, ymax=6,
  xtick={-5,...,3}, xticklabels={,$-4$,,$-2$,,,,$2$,},
  ytick={-10,-8,...,6}, yticklabels={,,,$-4$,,,$4$,},
  x=0.72cm, y=0.36cm]
\addplot[very thick, red, domain=-1.17:1.17, samples=80, <->] {-x^4-8};
\node[closed, red] at (axis cs:0,-8) {};
\node[font=\scriptsize, fill=white, inner sep=1pt, anchor=north east] at (axis cs:0,0) {$O$};
\node[font=\scriptsize, red, fill=white, inner sep=1pt] at (axis cs:-1.65,-9) {$g$};
\end{axis}
\end{tikzpicture}
\end{center}
\end{minipage}

\begin{minipage}[t]{0.55\textwidth}
\vspace{0pt}
\textbf{8. Error Analysis} Describe the error Terrence made in graphing the transformation of the cubic function $g(x)=x^3$ to $f(x)=-\tfrac12x^3+10$.
\lines[4]
\end{minipage}\hfill
\begin{minipage}[t]{0.42\textwidth}
\vspace{0pt}
\begin{center}
\begin{tikzpicture}
\begin{axis}[sheetgrid, xmin=-5, xmax=5, ymin=-25, ymax=25,
  xtick={-5,...,5}, xticklabels={,$-4$,,$-2$,,,,$2$,,$4$,},
  ytick={-25,-20,...,25}, yticklabels={,$-20$,,$-10$,,,,,,$20$,},
  x=0.5cm, y=0.1cm, title={Terrence's graph}, title style={font=\scriptsize, yshift=-4pt}]
\addplot[very thick, red, domain=-4.05:3.0, samples=80, <->] {0.5*x^3+10};
\node[font=\scriptsize, fill=white, inner sep=1pt, anchor=south east] at (axis cs:-0.1,10.6) {$10$};
\node[font=\scriptsize, fill=white, inner sep=1pt, anchor=north east] at (axis cs:0,0) {$O$};
\node[font=\scriptsize, red, fill=white, inner sep=1pt, anchor=east] at (axis cs:2.1,20) {$f(x)$};
\end{axis}
\end{tikzpicture}
\end{center}
\end{minipage}

\newpage
\wspart{Extension}{Practice 12 (DOK 3), when you finish}
\fbox{\parbox{\dimexpr\textwidth-2\fboxsep-2\fboxrule}{%
\small\raggedright An \textbf{even function} is symmetric about the $y$-axis: \mbox{for all $x$, $f(x)=f(-x)$.} An \textbf{odd function} is symmetric about the origin: \mbox{for all $x$, $f(-x)=-f(x)$.} The parent functions \mbox{$y=x^2, x^4, x^6,\ldots$} are even; \mbox{$y=x, x^3, x^5,\ldots$} are odd.}}

\textbf{Construct Arguments} Provide an example that demonstrates the following statement is not true.

\emph{If the degree of a function is an even number, then the function is an even function.}

Your function: $f(x)=\blank[4cm]$ \quad Its degree: \blank[1.2cm]

Show that it is not an even function (use $f(-x)$, or explain what your transformation does to the symmetry about the $y$-axis):
\lines[3]

\wspart{IXL}{Practice for the quiz}
Go to IXL, Algebra 2. Type the code in the search box. Do the starred skills first.

\vspace{2pt}
\begin{tabularx}{\textwidth}{l Y Y}
\toprule
\textbf{Code} & \textbf{Skill} & \textbf{Feeds} \\
\midrule
\ixlrow{KT8 $\bigstar$}{Describe function transformations}{Quiz 2 and 3: name each change from the parent function}
\ixlrow{RSN $\bigstar$}{Transformations of functions}{Quiz 2 and 4: the transformed graph and its equation}
\ixlrow{Q83}{Even and odd functions}{Optional: Quiz 1 (full Savvas quiz only) and the Extension}
\ixlrow{R7X}{Function transformation rules (from 1-2)}{Warm-up: $a$, $h$ and $k$ in $a\,f(x-h)+k$}
\ixlrow{XMK}{Match polynomial functions and graphs using end behavior (from 3-1)}{Warm-up: a negative leading coefficient flips the ends}
\bottomrule
\end{tabularx}

\vspace{4pt}
\textbf{SmartScore 70 or above on each starred ($\bigstar$) skill replaces your quiz score.}

\vspace{6pt}
\textbf{Before the quiz, finish these sentences.}

In $g(x)=4x^3+7$, the \textbf{parent function} is \blank[2cm], the \textbf{leading coefficient} 4 \blank[3.2cm], and adding 7 is a \textbf{translation} \blank[3.5cm].

When the leading coefficient is negative, the graph has a \textbf{reflection} across the \blank[2cm], so its ends \blank[4.5cm].

To get $f(x)=-x^4+3$ from $y=x^4$, first \blank[4cm], then \blank[4cm], because \blank[5cm].

\end{document}
